\subsection{A local certificate for all nine zero returns}
Let $P$ be the archived PA194CS register program, and let $P'$ be
PA194ZR.  A configuration consists of a source program counter and nine
natural-number registers $R_0,\ldots,R_8$.  The source instruction
$\operatorname{dec}(r,z,p)$ leaves a zero register unchanged and jumps to
$z$; on a positive register it subtracts one and jumps to $p$.

The programs have the same 3,646 instruction sites.  Their only differences
are the following nine replacements.  All site numbers are zero-based.
\begin{center}
\begin{tabular}{r|r|r}
site & old jump destination & new positive destination\\\hline
3 & 5 & 4\\
9 & 11 & 10\\
3526 & 3528 & 3527\\
3532 & 3535 & 3533\\
3544 & 3546 & 3545\\
3550 & 3553 & 3551\\
3553 & 3637 & 3554\\
3627 & 3629 & 3628\\
3633 & 3637 & 3634
\end{tabular}
\end{center}
At each listed site, $P'$ has
$\operatorname{dec}(3,\text{old destination},\text{positive destination})$.

Define the protected region
\[
 S=\{0,\ldots,9\}\ \cup\ \{3523,\ldots,3633\}
     \ \cup\ \{3638,\ldots,3645\}.
\]
It contains 129 sites, including all nine changed sites, and does not contain
the initial source counter 25.  Inspection of the complete control-flow
graphs of both $P$ and $P'$ gives exactly three edges entering $S$ from its
complement:
\[
 48\longrightarrow3638,\qquad
 49\longrightarrow3554,\qquad
 3519\longrightarrow3523.
\]
Each is the zero branch of a decrement of $R_3$.  No instruction at a site
in $S$ increments $R_3$.  The accompanying independent checker verifies
these facts by enumerating every edge of both complete instruction lists;
it uses no register-value bound, program simulation, or whole-program
positivity-mask certificate.

\paragraph{Local invariant.}
For either program, from source counter 25 with \emph{any} initial register
store, whenever the current counter lies in $S$, the current value of $R_3$
is zero.  This follows by induction on the number of source steps.  The
initial counter is outside $S$.  A step entering $S$ from outside is one
of the three zero branches above and therefore establishes $R_3=0$.
A step starting in $S$ cannot make a zero $R_3$ positive, because it does
not increment that register.  This includes a decrement of $R_3$: its
positive branch is then unavailable and its zero branch leaves the register
unchanged.

\paragraph{Exact source equivalence.}
Start $P$ and $P'$ from the same arbitrary register store at source counter
25.  After every number of source steps, their complete configurations
are equal.  Prove this by a second induction.  At every unchanged site
the instructions agree literally.  At a changed site the local invariant
gives $R_3=0$, so the new decrement leaves every register unchanged and
takes exactly the old jump's destination.  The two successors are thus
identical.  Consequently the source programs halt on exactly the same
initial stores, after the same number of source instructions.

This is a direct elementary argument about unbounded natural-number
registers.  Its finite graph premises and induction can be expressed in
PA.  It establishes the source-level transformation only: the separate
blank-tape initialization and universal physical transition contracts
remain necessary to transfer the conclusion to the displayed Turing
machine.  In particular, the rejected version using $R_6$ at the eight
shared returns had valid source contracts but failed blank-tape
initialization, so source equivalence alone was insufficient.
